%% kapa: exam=1 round=35 year=2024 period=2 form=A subject=accounting subject_ko=회계학 seq=18 number=58
%% source: Q-Net 감정평가사 기출문제 (cst00309) artlSeq=5213229 file=raw/1008/5213229/2교시 문제지 원본.pdf
%% conversion: kordoc 4.24.51 -> tools/gen1.py (auto); review=proofread-sonnet; flags=tab
\begin{question}{58}
(주)감평은 (주)대한이 발행한 사채(발행일 20×1년 1월 1일, 액면금액 ￦1,000, 표시이자율 연 8\%, 매년 말 이자지급, 20×4년 12월 31일에 일시상환)를 20×1년 1월 1일에 사채의 발행가액으로 취득하였다(취득 시 신용이 손상되어 있지 않음). (주)감평은 취득한 사채를 상각후원가로 측정하는 금융자산으로 분류하였으며, 사채발행시점의 유효이자율은 연 10\%이다. (주)감평은 (주)대한으로부터 20×1년도 이자 ￦80은 정상적으로 수취하였으나 20×1년 말에 상각후원가로 측정하는 금융자산의 신용이 손상되었다고 판단하였다. (주)감평은 채무불이행을 고려하여 20×2년부터 20×4년까지 현금흐름에 대해 매년말 수취할 이자는 ￦50, 만기에 수취할 원금은 ￦800으로 추정하였다. (주)감평의 20×1년도 포괄손익계산서의 당기순이익에 미치는 영향은? (단, 화폐금액은 소수점 첫째자리에서 반올림하며, 단수차이로 인한 오차는 가장 근사치를 선택한다.)
\par{\small\centering\begin{tabular}{|c|c|c|c|c|}\hline
\multirow{2}{*}{기간} & \multicolumn{2}{c|}{단일금액 ￦1의 현재가치} & \multicolumn{2}{c|}{정상연금 ￦1의 현재가치} \tabularnewline\cline{2-5}
 & 8\% & 10\% & 8\% & 10\% \tabularnewline\hline
3 & 0.7938 & 0.7513 & 2.5771 & 2.4868 \tabularnewline\hline
4 & 0.7350 & 0.6830 & 3.3120 & 3.1698 \tabularnewline\hline
\end{tabular}\par}
\begin{choices}
  \item ￦94 감소
  \item ￦94 증가
  \item ￦132 감소
  \item ￦226 감소
  \item ￦226 증가
\end{choices}
\end{question}
\answer{3}
